%% @Author: Ines Abdeljaoued Tej
%  @Date:   2018-06
%% @Class:  Graduation Project, ESSAI - Carthage University, Tunisia.

\section{Outil de Développement}

\begin{enumerate}
     \item \textbf{HyperText Markup Language}
    
 HTML permet la création de pages web basique. Il permet de mettre dans la zone la  Forme et le contenu d’une page web en utilisant une syntaxe composée.
\paragraph{\\}
 \renewcommand{\thefigure}{3}
    
    \begin{center}
        \begin{figure}[htbp]
    \centering
   \includegraphics[width=0.15\linewidth]{images/html (2).png} 
    \caption{Logo d'HTML}
   \paragraph{\\}\paragraph{\\}
    \label{fig:HyperText Markup Language}
\end{figure}
       
    \end{center}
    \paragraph{\\}

    \item \textbf{Framework Django}
    
 Django est une Plateform de développement web open-source en Python. Il fournit une structure ou architecture et des outils puissants pour la création rapide d'applications web    complexes et 100\% sécurisées.
     \renewcommand{\thefigure}{4}
    \begin{center}
        \begin{figure}[htbp]
    \centering
   \includegraphics[width=0.3\linewidth]{images/django1.png} 
    \caption{Logo de Django}
    \label{fig:Framework Django}
\end{figure}
        
    \end{center}

        
          \paragraph{\\}\paragraph{\\}\paragraph{\\}
     \item \textbf{Cascading Style Sheets}
     
 CSS est un langage de feuilles de style utilisé pour la présenter visuellement des documents HTML. Il permet un contrôle total de  l'apparence des éléments d'une page web, tels que la couleur, la police, la taille du texte, la mise en page, etc.

 \renewcommand{\thefigure}{5}
    
    \begin{center}
        \begin{figure}[htbp]
    \centering
   \includegraphics[width=0.1\linewidth]{images/cdd1.png} 
    \caption{Logo de Css}
    \label{fig:Cascading Style Sheets}
\end{figure}
        
    \end{center}
    \paragraph{\\}\paragraph{\\}\paragraph{\\}

      
     \item \textbf{Python}
    
Le langage Python est un langage de programmation polyvalent et puissant utilisé pour développer des logiciels système, des logiciels de bureau, des jeux vidéo et bien plus encore. Python fournit un contrôle précis sur la manipulation de la mémoire et permet la création de programmes performants et efficaces.

 \renewcommand{\thefigure}{6}
    
    \begin{center}
        \begin{figure}[htbp]
    \centering
   \includegraphics[width=0.1\linewidth]{images/python1.png} 
    \caption{Logo de Python}
    \label{fig:Python}
\end{figure}
        
    \end{center}
   
   \paragraph{\\}
   
   \paragraph{\\}
  
    \paragraph{\\}\paragraph{\\}
  
 \section{Premier pas vers l'UI de - "Safari"}

À partir de cette première interface utilisateur, L’utilisateur disposera de deux choix : se connecter ou s’inscrire. Il suffit donc d’entrer son nom d'utilisateur et son mot de passe pour avoir accès au menu autrement appèles accueil.

\renewcommand{\thefigure}{7}

\begin{center}
    \begin{figure}[htbp]
        \centering
        \includegraphics[width=\textwidth, height=0.8\textheight]{images/connexion_first.png} 
        \caption{Page de connexion}
        \label{fig:C++}
    \end{figure}
\end{center}

  \paragraph{\\}
  
    \paragraph{\\}
    \paragraph{\\}
  
  \paragraph{\\}
  \paragraph{\\}
  
    \paragraph{\\}
  
\section{Page D’inscription}

\renewcommand{\thefigure}{8}
  \paragraph{\\}
\begin{center}
    \begin{figure}[htbp]
        \centering
        \includegraphics[width=1\linewidth, height=0.65\textheight]{images/inscription_second.png} 
        \caption{Page d'inscription}
        \label{fig:Page d'inscription}
    \end{figure}
\end{center}
 
  
  
  \paragraph{\\}
  
    \paragraph{\\}
  
\section{Page de Publication de Blog}

\renewcommand{\thefigure}{9}
  \paragraph{\\}
\begin{center}
    \begin{figure}[htbp]
        \centering
        \includegraphics[width=1\linewidth, height=0.65\textheight]{images/publication_third.png} 
        \caption{Page de Publication de Blog}
        \label{fig:Page de Publication de Blog}
    \end{figure}
\end{center}
    
   
  \paragraph{\\}
  
    \paragraph{\\}
  
\section{Interface Admin}

\renewcommand{\thefigure}{10}
  \paragraph{\\}
\begin{center}
    \begin{figure}[htbp]
        \centering
        \includegraphics[width=1\linewidth, height=0.65\textheight]{images/Admin_Fourth.png} 
        \caption{Interface Admin}
        \label{fig:Interface Admin}
    \end{figure}
\end{center}
  
\end{enumerate}